\labsection{8}{Genetic Algorithms: Regression and Reproducibility}{sec:lab08-ga}

\begin{labproject}{Fit a medical-cost regression model with a Genetic Algorithm}
\textbf{Scenario.} You are testing whether a Genetic Algorithm can recover useful coefficients for a small regression problem and how its stochastic performance compares with an exact linear baseline.\\
\textbf{Software.} MATLAB.\\
\textbf{Goal.} Implement a real-valued GA, run it repeatedly on real data, and compare it with ordinary least squares.

\kagglesource{https://www.kaggle.com/datasets/mirichoi0218/insurance}{Medical Cost Personal Datasets}
\kaggleexample{https://www.kaggle.com/code/agarwalvishal00/medical-cost-personal-dataset-eda-and-regression}{Medical Cost Personal Dataset --- EDA and regression example}
\repofolder{lab08\_genetic\_algorithm/}

\begin{labmode}
Download \texttt{insurance.csv} manually from Kaggle and place it under \texttt{data/insurance/insurance.csv}. The file contains 1,338 rows and the fields age, sex, BMI, children, smoker, region, and charges.
\end{labmode}

\textbf{Required run.}
\begin{lstlisting}[style=mlmatlab]
matlab -batch "run('lab08_genetic_algorithm/src/main_lab08.m')"
\end{lstlisting}

\begin{mlsteps}
  \item \textbf{Verify the GA implementation.} Minimize $f(x)=(x-3)^2$ with a fixed seed and confirm that the best solution approaches $x=3$ before using real data.
  \item \textbf{Prepare the regression task.} Predict \texttt{charges} from age, BMI, number of children, and smoker status. Encode smoker as 0/1, standardize numeric predictors using training statistics, and use a fixed 80/20 train-test split with seed 42.
  \item \textbf{Define the chromosome.} Use five real-valued genes: intercept plus four coefficients. Fitness is training MSE. Use tournament or rank selection, crossover, mutation, and elitism; document all parameter values.
  \item \textbf{Run repeated trials.} Execute at least 10 GA runs with seeds 1--10. Save best training MSE, test RMSE, generations, and runtime for every run.
  \item \textbf{Compare with least squares.} Fit ordinary least squares on the identical training data and evaluate on the identical test data. Explain why least squares is expected to be preferable for this unconstrained linear model.
  \item \textbf{Sensitivity test.} Change one GA parameter---population size or mutation rate---and repeat the 10-run experiment. Compare median test RMSE and median runtime.
\end{mlsteps}

\begin{mltable}
\begin{tabularx}{\linewidth}{@{}>{\bfseries\leavevmode\color{MLNavy}}p{0.28\linewidth}XXX@{}}
\toprule
\rowcolor{MLPaleBlue!92}
Method & Test RMSE & Runtime & Notes \\
\midrule
GA --- median of 10 runs & & & \\
GA --- best run & & & \\
Least squares & & & \\
\bottomrule
\end{tabularx}
\end{mltable}

\requiredoutput{Saved split, 10-run results CSV, convergence plot, sensitivity comparison, and least-squares baseline.}
\end{labproject}

\begin{verification}
Use the same split and preprocessing for the GA and least-squares baseline. Do not report only the best stochastic run; the median and variability across seeds are part of the result.
\end{verification}

\begin{labdeliverable}
Push the completed lab and create tag \texttt{lab08-final}. The README must report median/standard-deviation test RMSE across seeds, the least-squares RMSE, the sensitivity result, and one justified conclusion about when a GA would be useful despite the baseline result.
\end{labdeliverable}
